\subsection{Подгруппы канонического теста}
\label{app:subgroup-details}
\begin{table}[!htbp]
\centering\small
\caption{Числа исходов замороженной коэффициентной KAN и обеих LUT (решения совпадают), а также число несертифицированных решений по интервалу, вычисленному от оценки LUT.}
\label{tab:subgroup-counts}
\begin{tabular}{lrrrrrr}
\toprule
Подгруппа & TN & FP & FN & TP & $U_{513}$ & $U_{1025}$\\
\midrule
Весь тест & 9786 & 214 & 896 & 31313 & 12 & 3\\
Совпадающие & 1609 & 25 & 0 & 3567 & 0 & 0\\
Без точного совпадения & 8177 & 189 & 896 & 27746 & 12 & 3\\
\bottomrule
\end{tabular}
\end{table}
Сохранённая вещественная KAN на восстановленных координатах до Q12 даёт F1/BA 0.983191/0.975946 на всём тесте, 0.996508/0.992350 на совпадающих строках и 0.981507/0.973684 на остальных. Для несовпадающей подгруппы её $(TN,FP,FN,TP)=(8176,190,857,27785)$. Этот расчёт условен относительно восстановления координат: исходный fitted-преобразователь не сохранён. Совпадение восстановленного Q12 на сохранённых строках не задаёт преобразователь нового потока. Полная точность всех метрик и определения трёх разных сертификатов сохранены в CSV; здесь используется именно условие от оценки LUT, как в табл.~\ref{tab:signed_certificate}.
